\documentclass[11pt]{amsart}
\usepackage{amsmath,amssymb}
\usepackage[hidelinks,pdfusetitle]{hyperref}
\allowdisplaybreaks[2]
\setlength{\emergencystretch}{3em}
\raggedbottom
\numberwithin{equation}{section}
\newtheorem{theorem}{Theorem}[section]
\newtheorem{proposition}[theorem]{Proposition}
\newtheorem{corollary}[theorem]{Corollary}
\newtheorem{lemma}[theorem]{Lemma}
\theoremstyle{remark}
\newtheorem*{remark}{Remark}

\title[M\"obius-weighted trilinear Kloosterman sums]{M\"obius-weighted trilinear Kloosterman sums at composite moduli and Goldbach remainder rectangles}
\hypersetup{pdftitle={M\"obius-weighted trilinear Kloosterman sums at composite moduli and Goldbach remainder rectangles}}
\author{KenRui Lin}
\address{Independent Researcher, Guangzhou, China}
\email{longinuslkr@gmail.com}
\subjclass[2020]{Primary 11L05, 11L07; Secondary 11P32, 11F72}
\keywords{Kloosterman fractions, M\"obius function, reciprocal energy, Goldbach problem, Poisson summation}
\date{}

\begin{document}
\begin{abstract}
We prove power savings over \(X^2\) for a trilinear sum with a M\"obius
weight, an essentially bounded second coefficient, and phases
\(e(Nh\bar a/b)\) with \(a,b\asymp X=N^{1-\theta}\) and
\(1\le J\le X^2/N\). A double M\"obius truncation, completion, and
averaged fourth and sixth reciprocal moments give a saving of
\(N^{(17\theta-8)/10-\varepsilon}\) for \(8/17<\theta\le1/2\).
Using a local spectral average of Pascadi in the short Type II range
gives a saving of \(N^{(15\theta-7)/10-\varepsilon}\) for
\(7/15<\theta\le1/2\). Both estimates are unconditional and uniform in
the prime factorization of \(N\). We compute the obstruction at
\(7/15\) for this combination of estimates and the feasible region
of the elementary fourth-moment argument in two variables.
Under the Riemann hypothesis, these estimates give a power saving
for smooth and sharp right-end dyadic rectangles of the Goldbach
long-divisor remainder.
\end{abstract}
\maketitle

\section{Introduction}

Let \(e(t)=e^{2\pi it}\) and let \(\bar a\) denote the inverse of \(a\)
modulo \(b\) in an expression with denominator \(b\). For \(N\ge1\) put
\(X=N^{1-\theta}\) and \(H=X^2/N\). Given a fixed smooth function
\(\Phi\) supported in \([1,2]^3\) and a complex sequence
\(\|\beta\|_\infty\ll_\varepsilon N^\varepsilon\) for every \(\varepsilon>0\),
write
\begin{equation}\label{eq:main-sum}
 S(X,J)=
 \sum_{\substack{a,b\asymp X\\(a,b)=1}}\mu(a)\beta_b
 \sum_{h\asymp J}\Phi(a/X,b/X,h/J)e(Nh\bar a/b).
\end{equation}
The support of \(\Phi\) specifies the dyadic ranges. All integers in
the sums are positive unless a frequency is explicitly signed.

\begin{theorem}[Elementary \(8/17\) bound]\label{thm:elementary}
For \(8/17<\theta\le1/2\) and \(1\le J\le H\),
\[
 |S(X,J)|\ll_{\Phi,\varepsilon} X^2 N^{-(17\theta-8)/10+\varepsilon}.
\]
The analytic inputs are Poisson summation and the full-modulus Weil
bound.
\end{theorem}

\begin{theorem}[Spectral \(7/15\) bound]\label{thm:spectral}
For \(7/15<\theta\le1/2\) and \(1\le J\le H\),
\[
 |S(X,J)|\ll_{\Phi,\varepsilon} X^2 N^{-(15\theta-7)/10+\varepsilon}.
\]
The estimate is unconditional and uniform for all \(N\), in particular
for arbitrary composite \(N\).
\end{theorem}

In both theorems the implied constants allow the constants in the
hypothesis on \(\|\beta\|_\infty\). The elementary argument using only
the fourth reciprocal moment gives
\[
 |S(X,J)|\ll_{\Phi,\varepsilon}X^2N^{-(19\theta-9)/12+\varepsilon}
 \qquad(9/19<\theta\le1/2).
\]

\begin{proposition}[Obstruction at \(7/15\)]\label{prop:obstruction}
Write \(X=N^x\), \(H=N^{j_0}\), and \(L=N^s\) for the short Type II
factor. At \(j_0=x/8\) and \(s=3j_0\), the critical exponents of
the reversed sixth and fourth moments both equal zero. The spectral
second-moment nonzero-frequency exponent equals \(x/64>0\).
Consequently the four-piece argument of this paper has exact
strict-saving threshold \(7/15\).
\end{proposition}

\begin{proposition}[Unbalanced feasible region]\label{prop:unbalanced}
Let \(U=N^\alpha,\ V=N^\beta\) with \(\alpha,\beta\ge0\),
\(p=\min(\alpha,\beta)\), \(q=\max(\alpha,\beta)\), and \(p+q\le1\).
At the nominal scales \(A=N/U,\ B=N/V,\ J\le AB/N\), allow either
orientation of the M\"obius-weighted variable and modulus in the
rectangle. The elementary nine-exponent method consisting of Type I,
the elementary second moment, and the fourth moment gives a strict
power saving over \(AB\) throughout the following feasible region:
\begin{equation}\label{eq:region}
\begin{aligned}
 &(2p+q>1\ \land\ 3p+5q>4)\\
 {}\lor{}&(4p+3q>3\ \land\ 8p+11q>9)\\
 {}\lor{}&(2p+q>1\ \land\ 2q-p>1).
\end{aligned}
\end{equation}
These inequalities describe the strict feasibility of the stated
method. This region strictly contains the one obtained by direct
insertion in Bettin--Chandee's Theorem~1,
\(19p+14q>16\) and \(8p+7q>7\).
\end{proposition}

For positive \(U,V\) define the sharp and smooth versions by
\begin{align}
 T_{U,V}(N)&=\sum_{\substack{au+bv=N\\ U\le u<2U\\ V\le v<2V}}
                    \mu(a)\mu(b), \label{eq:sharpT}\\
 T_{U,V}^{f,g}(N)&=\sum_{au+bv=N}\mu(a)\mu(b)f(u/U)g(v/V),
 \label{eq:smoothT}
\end{align}
where \(f,g\in C_c^\infty((0,\infty))\) are fixed and supported on
\([1,2]\). The same smooth estimate holds for fixed compact supports
inside \((0,\infty)\).

\begin{corollary}[Right-end Goldbach rectangles under RH]\label{cor:rectangles}
Assume the Riemann hypothesis for \(\zeta(s)\).
Let \(U,V\asymp N^\theta\) with fixed comparability constants.
For each fixed \(8/17<\theta<1\), the elementary method yields some
\(\eta=\eta(\theta)>0\) such that
\(T_{U,V}(N),T_{U,V}^{f,g}(N)\ll N^{1-\eta}\).
With the spectral input, the conclusion holds for \(7/15<\theta<1\).
The constants are uniform in the factorization of \(N\) and may
depend on \(\theta,f,g\) and the comparability constants.
\end{corollary}

The restriction \(\theta<1\) allows arbitrary fixed comparability
constants. At \(\theta=1\), taking \(U=V=9N/20\) forces \(a=b=1\)
and gives \(T_{U,V}(N)=N/10+O(1)\). For \(U=V=N^\theta\) with
\(\theta\ge1\), the equation has no solutions.

For the Goldbach application use the additive convolution
\[
 \mathcal B(F,G)(N)=\sum_{m=1}^{N-1}F(m)G(N-m).
\]
Convolving \(\Lambda*\mathbf1=\log\) with \(\mu\) gives
\(\Lambda(m)=\sum_{au=m}\mu(a)\log u\).
Consider a long-divisor component
\[
 R_Z(m)=\sum_{\substack{au=m\\a>Z}}\mu(a)w(u).
\]
It contributes to \(\mathcal B(R_Z,R_Z)(N)\) the same rectangle sums
with the additional weights \(w(u)w(v)\). For \(w(u)=\log u\), the
normalized factor \(\log(Ut)/\log(2N)\) has bounded derivatives on
the fixed support in \(t=u/U\). It can be included in the smooth
weights, or in the approximants used for the sharp rectangles.
These are the right-end rectangles in the long-divisor decomposition of the
GRH Goldbach problem. The estimates here cover the indicated right-end
scales. Control of the remaining intermediate scales is required for
a conclusion on the full Goldbach problem.

\subsection{Comparison with earlier estimates}
Bettin--Chandee's trilinear estimate allows three arbitrary
\(\ell^2\) coefficients and a separate fixed multiplier in the
phase \cite[Theorem~1]{BC}. Its direct balanced \(X^2\)-saving
threshold here is \(16/33\), obtained by putting \(p=q=\theta\)
in the two inequalities of Proposition~\ref{prop:unbalanced}.
Our extra structure is the M\"obius coefficient and the product
parameter \(Nh\). Bourgain--Garaev use rational reciprocal counts
and moment estimates to control Kloosterman sums
\cite{BGprime,BGrings}. Our application uses reciprocal energies
averaged over the modulus and divides the multiplier and modulus
by their common gcd before taking the moments.

Fouvry--Kowalski--Michel develop Type I and Type II estimates for
algebraic trace functions over prime fields and apply them to
prime and M\"obius weights \cite{FKM}. Here the modulus varies over
all integers in a dyadic range. The local spectral input is
Pascadi's \emph{Forum of Mathematics, Pi} paper
\cite[Corollary~5.5 and (5.20)]{PascadiPi}. Pascadi's separate
preprint gives arbitrary-coefficient estimates for Kloosterman
fractions \cite{PascadiFractions}. In its Corollary~4.8, substituting
coefficient lengths \(O(X)\), averaging length \(O(J)\) and fixed
multiplier \(N\) gives a term \(\sqrt{NJ}/X\) in the factor
multiplying the norm bound. At \(J\asymp H\) this term is of order
one, and the norm prefactor evaluated with the size bounds on our
coefficients is \(X^2J N^{o(1)}\). The direct substitution thus
does not yield a uniform power saving over \(X^2\).
The main small-numerator estimate of Dong--Robles--Zeindler
\cite[Theorem~1.1]{DRZ} has the size condition \(a\ll X^{1/10}\) on the numerator
of \(e(a\bar m/n)\). Here the numerator is \(Nh\), of order \(X^2\)
when \(J\asymp H\). The small parameter \(J\) by itself does not
verify their numerator condition. Shparlinski \cite{Shp} proves
bounds for trilinear sums of Kloosterman fractions in
parameter ranges where he improves the Bettin--Chandee bound. We record this
as related work and leave a numerical comparison with his full
parameter ranges open.

\section{The M\"obius decomposition and completion}
\label{sec:completion}

Write \(e(t)=\exp(2\pi i t)\), \(\widehat f(\xi)=\int_{\mathbb R}f(t)e(-t\xi)\,dt\), and \((a,b)=\gcd(a,b)\), with gcds taken positive. A bar denotes an inverse in the current modulus. All summation variables are positive integers unless otherwise specified. A scale \(D\) restricts its variable to a fixed compact interval of positive multiples of \(D\). Constants may depend on those intervals and on finitely many smooth seminorms. In intermediate estimates \(N^\varepsilon\) absorbs divisor bounds, the coefficient bounds, and fixed powers of \(\log(2N)\). One may use a smaller \(\varepsilon\) in each such step. The multiplier may be any integer \(n_0\) with \(1\le n_0\le N\).

For cutoffs \(R_1,R_2\ge1\), put \(u_i(m)=\mu(m)\mathbf1_{m\le R_i}\), \(v_i=\mu-u_i\), and let \(\mathbf1\) be the constant-one arithmetic function. Then
\begin{equation}\label{eq:mobius}
 \mu=u_1+u_2-u_1*u_2*\mathbf1+v_1*v_2*\mathbf1 .
\end{equation}
Indeed, \(\mu*\mathbf1=\delta_1\), since
\(\sum_{d\mid m}\mu(d)=\prod_{p\mid m}(1-1)\) for \(m>1\), with value
\(1\) for \(m=1\). Expanding \(v_1*v_2\), subtracting \(u_1*u_2\), and convolving
with \(\mathbf1\) gives \(\mu-u_1-u_2\). Take \(R_1=R_2=R\).
On the support \(a\asymp A\) with \(a>R\), the first two terms vanish. The
negative term becomes \(-\sum_{dm=a,\ d\le R^2}d_R(d)\), where
\(d_R(d)=\sum_{d_1d_2=d,\ d_i\le R}\mu(d_1)\mu(d_2)\) has absolute value
at most \(\tau(d)\). The positive term becomes
\(\sum_{\ell n=a,\ \ell,n>R}\mu(\ell)\sum_{v\mid n,\ v>R}\mu(v)\),
whose two coefficients are bounded by \(1,\tau(n)\).
These are the Type I and Type II pieces.

Choose a smooth \(\psi\) equal to \(1\) on \([0,1]\) and \(0\) on
\([2,\infty)\). The functions
\(\psi(t/2^{i+1})-\psi(t/2^i)\), \(i\in\mathbb Z\), sum to \(1\) for
\(t>0\) by telescoping. Use their products in the two factor variables.
Only \(O((\log(2N))^2)\) pairs meet the support, since
\(\ell n\asymp A\) and \(\ell,n>R\). In each pair rename the shorter
factor \(\ell\) and the longer \(n\), exchanging the coefficients when
needed. Their scales \(L,W\) satisfy \(L\gg R,\ L\ll A^{1/2}\), and
\(LW\asymp A\). Both \(>R\) conditions are retained in the coefficients.

For a smooth weight \(\Phi\) on scales \(A,B,J\) in \((a,b,h)\), the
split weight is a smooth function \(G(v_1,v_2,v_3,v_4)\) of
\((\ell/L,n/W,b/B,h/J)\). It is the dyadic factors times
\(\Phi((LW/A)v_1v_2,v_3,v_4)\). Extend it smoothly by zero outside fixed
compact intervals, using auxiliary cutoffs in each coordinate as needed.
Fourier inversion writes it as
\[
 \int_{\mathbb R^4}\widehat G(\mathbf t)\,
 e(t_1\ell/L)e(t_2 n/W)e(t_3 b/B)e(t_4 h/J)\,d\mathbf t .
\]
Integration by parts gives
\(|\widehat G(\mathbf t)|\le C_m(1+|\mathbf t|)^{-m}\) for \(m>4\).
Absorb each exponential into its own variable's coefficient, with
absolute values unchanged. Keep fixed interval supports for the
coefficients by inserting indicators equal to \(1\) on the support of
\(G\). This gives an absolutely convergent integral of blocks of the form
\begin{equation}\label{eq:separated}
 \sum_{\substack{\ell\asymp L,\ n\asymp W,\ b\asymp B,\ h\asymp J\\
                  (\ell n,b)=1}}
 \alpha_\ell\gamma_n\beta_b\nu_h\,e\!\left(\frac{n_0h\overline{\ell n}}b\right).
\end{equation}
The Fourier \(L^1\) cost is bounded by finitely many seminorms. When
Cauchy--Schwarz eliminates \(\gamma_n,\beta_b\), replace their support
indicators in the remaining nonnegative second moment by fixed smooth
nonnegative majorants. Their derivatives are independent of the
Fourier parameters.

For integers \(P,k\) and \(b\ge1\), define
\[
 \operatorname{Kl}(P,k;b)=\sum_{v\bmod b}^{*}e((P\bar v+kv)/b).
\]
The full-modulus Weil bound in the form valid for all \(b\) and all
integral arguments, including imprimitive ones, is
\begin{equation}\label{eq:weil}
 |\operatorname{Kl}(P,k;b)|\le \tau(b)b^{1/2}(P,k,b)^{1/2}
 \qquad\text{\cite[Chapter~11]{IK}}.
\end{equation}
For \(k=0\) inversion permutes the units. Expanding their indicator by
M\"obius inversion gives \(c_b(P)=\sum_{d\mid(b,P)}d\mu(b/d)\).
For \(P\ne0\) we use
\begin{equation}\label{eq:ram-average}
 |c_b(P)|\le(b,P),\qquad
 \sum_{b\asymp B}(P,b)\ll B\tau(|P|).
\end{equation}
The formula for \(c_b\) is multiplicative in \(b\) and at \(p^\nu\)
gives \(\varphi(p^\nu),-p^{\nu-1},0\) according as \(v_p(P)\) is at
least \(\nu\), equals \(\nu-1\), or is smaller. For the average,
expand \((P,b)=\sum_{d\mid(P,b)}\varphi(d)\), obtained by grouping
\(1,\ldots,(P,b)\) by their gcd with \((P,b)\). Only \(d\ll B\) arise
and their multiples number \(O(B/d)\).

For every \(K>0\) and \(M>1\),
\begin{equation}\label{eq:weil-group}
 \sum_{k\ne0}(P,k,b)^{1/2}(1+|k|/K)^{-M}
 \ll_M K\sum_{d\mid b}d^{-1/2}.
\end{equation}
Bound the gcd square root by \(\sum_{d\mid(b,k)}d^{1/2}\). For
\(k=dt\ne0\), comparison of \((1+dx/K)^{-M}\) with its integral on
\(x\ge0\) bounds the sum over \(t\ge1\) by \(K/(d(M-1))\). This
includes \(b\mid P\). If \(K<1\), one can multiply the right side by
\(K^E\) for any fixed \(E>0\) by increasing the decay exponent on the
left by \(E\), since \(1+|k|/K\ge1/K\).

\subsection{Type I}
For fixed \(d,b,h\), the Type I sum vanishes if \((d,b)>1\).
Otherwise \(m\) is on scale \(Y=A/d\), with smooth weight \(f(m/Y)\)
and phase \(e(n_0h\bar d\,\bar m/b)\).
Poisson summation in each residue class modulo \(b\) gives
\[
 \frac{Y}{b}\sum_{k\in\mathbb Z}
   \widehat f(kY/b)\operatorname{Kl}(n_0h\bar d,k;b).
\]
The residue \(n_0h\bar d\) has gcd \((n_0h,b)\) with \(b\).
The zero frequency is \(\ll_\varepsilon YN^\varepsilon\) upon summing
over \(b\) by \eqref{eq:ram-average}. Summation over \(h,d\) gives
\(\ll_\varepsilon JA N^\varepsilon\), since
\(\sum_{d\le R^2}d^{-1}\ll\log(2N)\).
At nonzero frequencies, \eqref{eq:weil}--\eqref{eq:weil-group}
with \(K=B/Y\) give \(\ll_\varepsilon B^{1/2}N^\varepsilon\)
per \(b\), since \(Y/b\) cancels \(K\) up to bounded factors.
Summing gives \(\ll_\varepsilon JR^2 B^{3/2}N^\varepsilon\).
If all \(Y\gg BN^\delta\) for a fixed \(\delta>0\), the improved
decay following \eqref{eq:weil-group} makes this part smaller than any
fixed negative power of \(N\), by using sufficiently many derivatives.

\subsection{The elementary Type II second moment}
\label{sec:second}
Apply Cauchy--Schwarz on \((b,n)\) in \eqref{eq:separated}. Its
coefficient square sum is \(\ll_\varepsilon BW N^\varepsilon\).
The remaining second moment, with smooth nonnegative majorants
\(w_b(b/B),w_n(n/W)\), expands over \(h_1,h_2,\ell_1,\ell_2\).
Put \(D=\ell_1\ell_2\) and \(\Delta=h_1\ell_2-h_2\ell_1\).
Then \((D,b)=1\) and the phase is \(e(n_0\Delta\overline D\,\bar n/b)\).
When \(\Delta=0\) the number of tuples is \(\ll_\varepsilon JL N^\varepsilon\):
for each of the \(O(JL)\) choices of \((h_1,\ell_2)\), the number of
factorizations \(h_2\ell_1=h_1\ell_2\) is at most \(\tau(h_1\ell_2)\).
The diagonal moment is thus \(\ll_\varepsilon BWJL N^\varepsilon\).
All coefficient products cost at most \(N^\varepsilon\).

For \(\Delta\ne0\), completion of the \(n\) sum gives, per tuple,
\begin{equation}\label{eq:second-poisson}
 \sum_{\substack{b\asymp B\\(b,D)=1}}
 w_b(b/B)\frac{W}{b}
 \sum_{k\in\mathbb Z}\widehat w_n(kW/b)
 \operatorname{Kl}(n_0\Delta\overline D,k;b).
\end{equation}
For \(k=0\) the sum over \(b\) is \(\ll_\varepsilon WN^\varepsilon\)
by \eqref{eq:ram-average}, since \(n_0\Delta\ne0\).
For \(k\ne0\) it is \(\ll_\varepsilon B^{3/2}N^\varepsilon\)
by \eqref{eq:weil-group} with \(K=B/W\), including \(W>B\).
If \(W\gg BN^\delta\) these frequencies have arbitrarily strong power decay.
There are \(O(J^2L^2)\) tuples in each off-diagonal piece. Multiplying
each bound by the outer \(BW N^\varepsilon\) and taking square roots
gives respective ratios to \(AB\) bounded by \(N^\varepsilon\) times
\[
 (J/L)^{1/2},\qquad J/B^{1/2},\qquad
 J L^{1/2}B^{1/4}/A^{1/2}.
\]
This uses \(LW\asymp A\) in each normalization.

\subsection{The spectral estimate for the short Type II band}
\label{sec:spectral-completion}

We use the Kloosterman sum convention
\[
 S(m,n;b)=\sum_{v\bmod b}^{*}e\!\left(\frac{m\overline v+nv}{b}\right).
\]
In the completed second moment, let the two terms have indices
\((h_1,\ell_1)\) and \((h_2,\ell_2)\), and put
\[
 D=\ell_1\ell_2,\qquad
 \Delta=h_1\ell_2-h_2\ell_1.
\]
The phase after expansion is \(e(P\Delta\overline D\,\overline n/b)\),
where \(P=N\) on the balanced box. The proof below also allows any
positive integer \(P\le N\) for the numerator. If the signed frequencies
are negative, we use their absolute values for their scale and keep
their sign in \(\Delta\). All four variables are nonzero. After Poisson
in \(n\) with a smooth cutoff \(V(n/W)\) and a smooth modulus cutoff
\(U(b/B)\), the term for this quadruple and \(k\) is
\[
 \sum_{(b,D)=1}\frac{W}{b}U(b/B)\widehat V(kW/b)
           S(P\Delta\overline D,k;b),
 \qquad
 \widehat V(\zeta)=\int_{\mathbb R}V(v)e(-v\zeta)\,dv.
\]
Here \(U,V\) are fixed smooth cutoffs supported on compact subsets
of \((0,\infty)\). As in the elementary second moment, the Cauchy--Schwarz
factor outside this moment is \(BW\), and the arithmetic coefficients
on the quadruples are \(O_\varepsilon(N^\varepsilon)\).
We treat \(\Delta\ne0\), \(k\ne0\), and write \(m=|P\Delta|\) and
\(G=(D,m)\). Negating both arguments of a classical Kloosterman sum
leaves it unchanged by the substitution \(v\mapsto -v\).
Thus for a fixed sign of \(k\) we can use \(m\) as the positive argument,
with the corresponding consistent sign on the second argument.

\paragraph{The cusps and the completed weight.}
In \(\Gamma_0(D)\) take the cusps \(\infty,1=1/1\) and the matrices
\[
 \sigma_\infty=I,\qquad
 \sigma_1=\begin{pmatrix}\sqrt D&-1/\sqrt D\\
                           \sqrt D&0\end{pmatrix}.
\]
These agree with the normalization in \cite[(3.9)]{PascadiPi}
for \(D>1\), and for \(D=1\) the second matrix differs by an integral
translation from that normalization. The cusp \(1\) has width \(D\)
before scaling, since conjugation of a translation by
\(\bigl(\begin{smallmatrix}1&0\\1&1\end{smallmatrix}\bigr)\)
makes its lower left entry the negative of the translation parameter.
Both cusps have \(\mu(\mathfrak a)=1/D\) in Pascadi's notation. They
are the Atkin--Lehner cusps corresponding to \(s=D\) and \(s=1\),
respectively, in the parametrization \(rs=D\), \((r,s)=1\).

For \(\gamma=\bigl(\begin{smallmatrix}a_0&b_0\\D c_0&d_0\end{smallmatrix}\bigr)
\in\Gamma_0(D)\), the lower left entry of \(\gamma\sigma_1\) is
\(b\sqrt D\), where \(b=Dc_0+d_0\), and its lower right entry is
\(u\sqrt D\), where \(u=-c_0\). The determinant implies
\[
 (b,D)=1,\qquad
 (a_0+b_0)d_0-b_0 b=1,\qquad d_0=b+Du.
\]
Thus \(u\) is a unit modulo \(b\) and
\(a_0+b_0\equiv\overline D\,\overline u\pmod b\).
Conversely, if \(b>0\), \((b,D)=1\), and \(u\) is a unit modulo \(b\),
choose a representative for \(u\), put \(c_0=-u\), \(d_0=b+Du\),
take \(v\) with \(v d_0\equiv1\pmod b\), and set
\(b_0=(v d_0-1)/b\) and \(a_0=v-b_0\). These entries give a matrix of
determinant \(1\) in \(\Gamma_0(D)\). The lower right classes modulo
\(b\sqrt D\) are therefore exactly the classes \(u\sqrt D\) just
described, and each determines its upper left class. The definition of
the generalized sum gives
\[
 S_{\infty,1}(m,k;b\sqrt D)=S(m\overline D,k;b)
 \quad ((b,D)=1),
\]
which is the case \(r=D,s=1\) of \cite[(3.12)]{PascadiPi}.
The modulus scale in that formula is \(C=B\sqrt D\).

Partition the positive integers \(|k|\) into \(K<|k|\le2K\)
for dyadic \(K\ge1/2\), as in the \(n\)-sum of Pascadi's
corollary. Choose a smooth \(\chi\) supported in a fixed compact subset
of \((0,\infty)\) and equal to one on \([1,2]\). On each piece put
\[
 \begin{gathered}
 z=\frac{|k|}{K},\qquad y=\frac{b}{B}=\frac{c}{C},\qquad
 t=\frac{KW}{B},\\
 F_{t,\pm}(z,y)=\chi(z)\frac{U(y)}{y}\widehat V(\pm tz/y),
 \end{gathered}
\]
where the sign in \(F_{t,\pm}\) is the original sign of \(k\).
The weight in the spectral estimate is \(F_{t,\pm}\), with the
constant \(W/B\) outside. Each \(\partial_z^a\partial_y^b F_{t,\pm}\)
is a finite sum of bounded smooth factors on a fixed compact set times
\(t^i\widehat V^{(i)}(\pm tz/y)\), with \(i\le a+b\).
The transform \(\widehat V^{(i)}\) is the Fourier transform of
\((-2\pi i v)^i V(v)\), so \(M+i\) integrations by parts bound
each such product by \(O_{a,b,M}((1+t)^{-M})\).
Extending by zero and taking the Fourier transform in \(z\) gives
\[
 \widetilde F_{t,\pm}(\xi,y)=
    \int_{\mathbb R}F_{t,\pm}(z,y)e(-\xi z)\,dz,\qquad
 F_{t,\pm}(z,y)=\int_{\mathbb R}
    \widetilde F_{t,\pm}(\xi,y)e(\xi z)\,d\xi.
\]
Integrating by parts in \(z\) also gives, for all nonnegative integers
\(b,M,Q\),
\begin{equation}
 |\partial_y^b\widetilde F_{t,\pm}(\xi,y)|
   \ll_{b,M,Q}(1+t)^{-M}(1+|\xi|)^{-Q}.
 \label{eq:spectral-schwartz}
\end{equation}
For the Bessel form, set
\[\Psi_t(z,u)=\sqrt z\,F_{t,\pm}(z,\sqrt z/u),
  \qquad U_{\rm Bes}=C/\sqrt{mK}.\]
Let \(\widetilde\Psi_t\) be the Fourier transform in \(z\).
Fourier inversion gives
\[
 \begin{aligned}
 F_{t,\pm}(|k|/K,c/C)
 &=\frac Cc\int_{\mathbb R} e(\xi |k|/K)
 \frac{\widetilde\Psi_t(\xi,v/(4\pi/U_{\rm Bes}))}
      {v/(4\pi/U_{\rm Bes})}\,d\xi,\\
 v&=\frac{4\pi\sqrt{m|k|}}c .
 \end{aligned}
\]
The support of \(\Psi_t\) has \(z,u\) in fixed compact subsets of
\((0,\infty)\), so the same integration by parts bounds its Fourier
transform and every fixed \(u\)-derivative. This supplies the
Fourier separation used in \cite[(5.21)--(5.26)]{PascadiPi} and
preserves rapid decay in \(t\).

\paragraph{The spectral input and its parameters.}
We invoke \cite[Corollary 5.5, (5.20)]{PascadiPi}, numbered
Corollary 15 in arXiv:2404.04239v3. In its notation, \(q\) is the
level, \(\mathcal N\) the length of the positive spectral argument,
\(C\) the modulus scale, and \(Z_{\rm aux}\) the auxiliary smoothness
parameter. For a smooth weight \(\mathcal F\) supported with both
arguments of order \(1\), with
\(\partial_z^i\partial_y^j\mathcal F\ll_{i,j,\varepsilon}
Z_{\rm aux}^{i\varepsilon}\), the full estimate is
\begin{align}
 &\left|
 \sum_{\mathcal N<n\le2\mathcal N}a_n
 \sum_{c\in\mathcal C_{\mathrm{ab}}}
 \mathcal F(n/\mathcal N,c/C)S_{\mathrm{ab}}(m,\pm n;c)
 \right| \notag\\
 &\quad\ll_\varepsilon (qm\mathcal N C Z_{\rm aux})^{O(\varepsilon)}
 (1+\mathcal T)^{2\vartheta(q)}
 \frac{C^2 A_{\mathcal N}}{C+\sqrt{m\mathcal N}}\notag\\
 &\qquad{}\times
 \left(1+\frac{m\mathcal N}{C^2}+
       \frac{\sqrt{(q,m)m}}q\right)^{1/2}
 \left(1+\frac{m\mathcal N}{C^2}+
       \frac{\mathcal N}q\right)^{1/2}.
 \label{eq:pascadi-completion}
\end{align}
The cusps have \(\mu(\mathrm a)=\mu(\mathrm b)=1/q\), with
the scaling above. The parameters for \(a_n=e(n\alpha)\) are
\(A_{\mathcal N}=\sqrt{\mathcal N}\), \(Z_{\rm aux}=1\), and
\(Y_{\mathcal N}=\mathcal N/T_{\mathcal N}(\alpha)\) by (5.18), where
\(T_{\mathcal N}(\alpha)=
\min_{v\ge1,\ v\in\mathbb Z}(v+2\mathcal N\|v\alpha\|)\).
In particular \(T_K(0)=1\), and inserting \(v=1\) shows that
\(T_K(\xi/K)\le1+2|\xi|\).
In the normalization of (5.20), (5.25), and (5.27) one has
\[
 \begin{aligned}
 \mathcal T &=
 \frac{C}{\sqrt{\max(m,q^2/(q,m))\,\max(\mathcal N,q)\,
                        Y_{\mathcal N}}}
   =\frac{C/\sqrt{m\mathcal N}}{\sqrt{X_1X_2(0)}},\\
 X_1&=\max\!\left(1,\frac{q^2}{(q,m)m}\right),\qquad
 X_2(0)=\max(1,q/\mathcal N)Y_{\mathcal N}.
 \end{aligned}
\]
In particular, writing \(\mathcal T=\mathcal T_0/\sqrt{Y_{\mathcal N}}\)
gives \(\mathcal T_0\le C\sqrt{(q,m)}/q^{3/2}\).
The quantity \(\vartheta(q)\) is the maximum exceptional spectral
parameter at level \(q\), with zero used for an empty exceptional
spectrum, and \(\vartheta(q)\le7/64\) by the Kim--Sarnak bound
\cite[Theorem 3.3]{PascadiPi}.

We set \(\mathcal N=K\), \(q=D\), and take \(a_n=1\) when (5.20)
is applied directly with weight \(F_{t,\pm}\), normalized using the
\((1+t)^{-M}\) bounds. Then \(X_2(0)=\max(K,D)\) exactly and
\begin{equation}
 \mathcal T\le\frac{B\sqrt D}{\sqrt{mK}\sqrt{\max(K,D)}}
              \le\frac{B}{\sqrt{mK}}.
 \label{eq:spectral-level-cancellation}
\end{equation}
The same bound after integration follows from the explicit separated
form. Insert a fixed cutoff in \(z\) equal to one on the support of
\(\chi\) and use the first Fourier integral with
\(\alpha=\xi/K\). In this case
\(X_2(0)=\max(K,D)/T_K(\xi/K)\); the resulting factor
\((1+2|\xi|)^{\vartheta(D)}\) in the bound is integrable against
\eqref{eq:spectral-schwartz}. In the Bessel version the perturbed
large-sieve parameter in (5.25) is
\(\max(K,D)/(1+|\xi|^2)\) for the zero base twist, which is
controlled in the same way by rapid Fourier decay. If this displayed
value falls below \(1\), one may use the unweighted spectral large
sieve at \(1\).

\paragraph{Length, exceptional factor and the completed sum.}
First suppose \(B=X\), \(LW=X\), and \(P=N\).
Write \(X=N^x\), \(L=N^s\), \(H=X^2/N=N^{j_0}\), and assume
\(J\le H\), \(j_0\le x/8\), \(j_0\le s\le x/2\).
All the scales in the short spectral band meet these inequalities.
We have \(D\asymp L^2\), \(N\le m\ll NJL\), and \(K=tL\).
Consequently
\[
 \frac{mK}{C^2}\ll t\frac{NJ}{X^2}\le t,\qquad
 \frac KD\ll t/L\le t,\qquad
 \mathcal T\le (H/K)^{1/2}.
\]
In \eqref{eq:pascadi-completion} use
\(C^2/(C+\sqrt{mK})\le C\), \(A_K=\sqrt K\), and
\(\sqrt{1+a+b}\le 1+\sqrt a+\sqrt b\) for \(a,b\ge0\).
After enlarging \(M\) in \eqref{eq:spectral-schwartz}, either sign
of the \(K\)-piece, including the factor \(W/X\), is
\begin{equation}
 \ll_{\varepsilon,M} N^\varepsilon(1+t)^{-M}
 W L \sqrt K(1+\mathcal T)^{2\vartheta(D)}
 \left(1+\frac{G^{1/4}m^{1/4}}{L}\right).
 \label{eq:spectral-piece}
\end{equation}
For \(0\le\lambda=2\vartheta(D)\le7/32\), we have
\(\sqrt K(1+\mathcal T)^\lambda\ll
\sqrt K+K^{(1-\lambda)/2}H^{\lambda/2}\).
Since \(H\le L\), this is
\(\ll\sqrt L(1+t)^{1/2}\). The rapid decay for \(K>L\)
and the \(O(\log(2N))\) blocks \(K\le L\) give a total at most
\begin{equation}
 N^\varepsilon W L^{3/2}
       \left(1+G^{1/4}X^{1/2}L^{-3/4}\right)
 \ll_\varepsilon N^\varepsilon W X^{1/2}L^{3/4}G^{1/4}.
 \label{eq:spectral-fixed-tuple}
\end{equation}
Here \(m\ll NJL\le X^2 L\) and \(X^{1/2}L^{-3/4}\ge X^{1/8}\ge1\).
To justify the \(N^\varepsilon\) factor uniformly over the unbounded
dyadic sum, use \((D m K C)^{O(\varepsilon)}
\ll N^{O(\varepsilon)}(1+t)^{O(\varepsilon)}\) before choosing the
decay order and then use a smaller input \(\varepsilon\).

The same row is valid even if the larger parameter
\(U_{\rm Bes}=C/\sqrt{mK}\) is retained in
\eqref{eq:spectral-piece}. Indeed, the term there with
\(G^{1/4}m^{1/4}/L\) leads, apart from its factors independent
of \(m,K\), to
\[
 \sqrt K m^{1/4}
    +(XL)^\lambda K^{(1-\lambda)/2}m^{1/4-\lambda/2}.
\]
Both exponents on \(m\) are positive. Substitute \(m\ll NJL\);
upon writing \(K=tL\), the second term is bounded by a constant
times
\[
 (1+t)^{1/2}\sqrt L\,(NJ L)^{1/4}(H/J)^{\lambda/2}
 \le (1+t)^{1/2}\sqrt L\,(X^2 L)^{1/4},
\]
because \((J/H)^{1/4-\lambda/2}\le1\).
The first term has the same bound without \( (H/J)^{\lambda/2}\).
For the term \(1\) in \eqref{eq:spectral-piece}, use \(m\ge N\)
to get
\[
 \sqrt K(1+U_{\rm Bes})^\lambda
 \ll \sqrt L(1+t)^{1/2}\bigl(1+(HL)^{\lambda/2}\bigr).
\]
The inequality
\[
 \frac{\lambda}{2}(j_0+s)\le
 \frac{7}{64}\left(\frac x8+\frac x2\right)
 <\frac x8\le\frac x2-\frac{3s}{4}
\]
holds for \(x>0\) and dominates this last loss by
\(X^{1/2}L^{-3/4}\). The bounds include \(K=1\), \(J=1\), and
the endpoint \(j_0=x/8\).

\paragraph{The average over quadruples.}
For every \(\Delta\ne0\), and for every \(P\) including integers
with many small prime factors,
\begin{align*}
 (D,P\Delta)^{1/4}
 &\le(\ell_1,P)^{1/4}(\ell_2,P)^{1/4}
       (\ell_1,\Delta)^{1/4}(\ell_2,\Delta)^{1/4}\\
 &\le(\ell_1,P)^{1/4}(\ell_2,P)^{1/4}
       (\ell_1,\ell_2)^{1/2}
       (\ell_1,h_1)^{1/4}(\ell_2,h_2)^{1/4}.
\end{align*}
The first inequality uses \((uv,w)\mid(u,w)(v,w)\) and
\((u,vw)\mid(u,v)(u,w)\), both verified by taking the valuation
at each prime. For the second, reduce \(\Delta\) modulo each
\(\ell_i\) and use the second divisibility again. For each \(\ell\),
\[
 \sum_{1\le |h|\ll J}(\ell,h)^{1/4}
 \ll\sum_{\substack{e\mid\ell\\ e\ll J}} e^{1/4}\frac{J}{e}
 \ll J\tau(\ell).
\]
To see the first inequality, replace the greatest common divisor
to the \(1/4\) power by the sum of the \(1/4\) powers of its divisors
and count nonzero multiples of \(e\).
Likewise,
\[
 \begin{aligned}
 \sum_{\ell_1,\ell_2\asymp L}(\ell_1,\ell_2)
 &\ll L^2\sum_{e\ll L}\frac1e\ll L^2\log(2L),\\
 \sum_{\ell\asymp L}(\ell,P)^{1/2}
 &\ll L\sum_{d\mid P}d^{-1/2}\le L\tau(P).
 \end{aligned}
\]
In the first sum we majorize a greatest common divisor by the sum
of its divisors and count pairs of multiples. In the second we
majorize its square root by the sum of the square roots of its divisors
and count multiples. Cauchy--Schwarz on the \((\ell_1,\ell_2)\) sum
now gives
\begin{equation}
 \sum_{\substack{\ell_1,\ell_2\asymp L,\ |h_i|\asymp J\\\Delta\ne0}}
 (D,P\Delta)^{1/4}
 \ll_\varepsilon J^2 L^2 N^\varepsilon.
 \label{eq:spectral-gcd-average}
\end{equation}
In particular the bound uses only divisor bounds and applies to
primorial and factorial \(P\) with the same constants.

Using \eqref{eq:spectral-fixed-tuple} and
\eqref{eq:spectral-gcd-average} gives a bound
\(N^\varepsilon J^2 W X^{1/2}L^{11/4}\) for the portion of
the completed second moment under consideration. Multiplying by
the outside factor \(XW\) and taking a square root gives
\[
 N^\varepsilon J W X^{3/4}L^{11/8}
   = N^\varepsilon J X^{7/4}L^{3/8}.
\]
After division by \(X^2\) the exponent is
\(j+3s/8-x/4\). The term \(1\) in
\eqref{eq:spectral-piece}, with the precise exceptional parameter,
separately gives \(N^\varepsilon J^2WL^{7/2}\) before the outside
factor and square root. Its normalized exponent is
\(j+3s/4-x/2\), which is at most \(j+3s/8-x/4\), since
\(3s/8-x/4\le-x/16\). The diagonal and the \(k=0\) terms
retain their elementary estimates.

\paragraph{Stability at the Poisson scales for rectangles.}
For use with the outer Poisson summation, suppose
\(Z=N^{x'}\), \(L W=A\), \(L\ll\sqrt Z\),
and \(ZN^{-\rho}\ll A,B\ll Z\). Assume
\(J\le N^c\), \(Z^2/P\le N^c\), and \(0\le c\le x'/8\), and
write \(s=\log_N L\ge c\). The natural frequency length is
\(K_0=B/W=BL/A\), so
\(L N^{-\rho}\ll K_0\ll L N^\rho\), and \(K=tK_0\).
The estimates for the two length ratios become
\[
 \frac{mK}{C^2}\ll t\,\frac{PJ}{AB},\qquad
 \frac KD\ll t\,N^\rho/L .
\]
In the application \(P=N/d_0\), \(Z=X/d_0\),
and \(J\ll (Z^2/P)N^\delta\). Thus the right sides are
\(\ll t N^{2\rho+\delta}\) and \(\ll tN^\rho\), respectively.
Also \(C\asymp BL\),
\(\mathcal T\le B/\sqrt{mK}\ll(N^c/K)^{1/2}\), and
\(m\ll Z^2 L N^\delta\). The argument bounding
\(\sqrt K(1+\mathcal T)^\lambda\) now gives
\(\ll\sqrt L N^{\rho/2}(1+t)^{1/2}\).
These substitutions in the full two factors of
\eqref{eq:pascadi-completion} cost \(N^{O(\rho+\delta)}\).
In particular the bound for an individual quadruple is
\(N^{O(\rho+\delta)+\varepsilon}W Z^{1/2}L^{3/4}G^{1/4}\).
Equation \eqref{eq:spectral-gcd-average} still applies. The outside
factor is \(BW\), hence the bound after taking the square root is
\begin{equation}
 \ll_\varepsilon N^{O(\rho+\delta)+\varepsilon}
             J Z^{7/4}L^{3/8},
 \label{eq:spectral-rectangular}
\end{equation}
using \(B\ll Z\) and \(W=A/L\ll Z/L\).

\section{Averaged reciprocal energies and the higher-moment bounds}
\label{sec:higher-moments}

The rational reciprocal-energy estimate is due to Karatsuba and is
recorded by Bourgain and Garaev \cite[Lemma~4]{BGprime}. We first give the
least-common-multiple argument and then average over the actual moduli.
For comparison, Bourgain and Garaev \cite{BGrings} study reciprocal
energies and multilinear Kloosterman sums at a given general modulus.

\begin{lemma}[Karatsuba's least-common-multiple bound]
\label{lem:target-lcm}
Let \(f_1,\ldots,f_m\) be positive integers, let
\(\sigma_1,\ldots,\sigma_m\in\{-1,1\}\), and set
\(M=\operatorname{lcm}(f_1,\ldots,f_m)\).
If \(\sum_{i=1}^m \sigma_i/f_i=0\), then
\[
 M^2\ \bigm|\ \prod_{i=1}^m f_i .
\]
\end{lemma}
\begin{proof}
Multiplication by \(M\) gives \(\sum_i \sigma_i(M/f_i)=0\) in
\(\mathbb Z\).
Fix a prime \(p\mid M\) and write \(v=v_p(M)=\max_i v_p(f_i)\).
If exactly one index attains \(v\), then exactly one summand of this
integer relation is a unit modulo \(p\), which is impossible.
Thus at least two indices attain \(v\), and
\(v_p(\prod_i f_i)=\sum_i v_p(f_i)\ge 2v\).
This holds for every prime dividing \(M\), proving the divisibility.
\end{proof}

For \(k\in\{2,3\}\), let \(E_{q,k}(T)\) count ordered tuples
\(f_1,\ldots,f_{2k}\), all of size \(T\) and coprime to \(q\), with
\[
 \bar f_1+\cdots+\bar f_k
    \equiv \bar f_{k+1}+\cdots+\bar f_{2k}\pmod q .
\]
The intervals defining size \(T\) lie in a fixed positive multiple of
\([T,2T]\). All lengths in this proof are at most a fixed power of \(N\).
We will use \(\tau(m)\ll_\epsilon m^\epsilon\) for positive integers
\(m\). Indeed, in the prime factorization of \(m\), for
\(p\ge 2^{1/\epsilon}\) and \(a\ge0\) one has \(a+1\le 2^a\le
p^{\epsilon a}\); for each of the finitely many smaller primes the
supremum of \((a+1)p^{-\epsilon a}\) is finite. By choosing the input
\(\epsilon\) small enough, each fixed product of divisor bounds costs
\(O_\varepsilon(N^\varepsilon)\).

For a tuple satisfying the corresponding equality over \(\mathbb Q\),
Lemma~\ref{lem:target-lcm} gives
\(M\le(\prod_i f_i)^{1/2}\ll_k T^k\).
Each \(f_i\) divides \(M\), so the total number of such tuples is bounded
by
\[
 \sum_{M\ll_k T^k}\tau(M)^{2k}\ll_{k,\varepsilon}T^kN^\varepsilon .
\]
There are \(O(Q)\) moduli \(q\asymp Q\) for \(Q\gg 1\), so these tuples
contribute \(O_{k,\varepsilon}(QT^kN^\varepsilon)\).
For each remaining tuple, form the nonzero integer
\[
 F=\sum_{i=1}^{2k}\sigma_i\prod_{r\ne i} f_r,\qquad
 \sigma_i=
 \begin{cases}
   1&i\le k,\\
  -1&i>k .
 \end{cases}
\]
Multiplication of the congruence by \(\prod_i f_i\) shows that
\(q\mid F\), with \(\lvert F\rvert\ll_k T^{2k-1}\).
There are \(O_\varepsilon(N^\varepsilon)\) possible \(q\) for any such
tuple, and \(O(T^{2k})\) tuples. Consequently
\begin{equation}
 \sum_{q\asymp Q}E_{q,k}(T)
    \ll_{k,\varepsilon}N^\varepsilon(QT^k+T^{2k}).
 \label{eq:target-mean-energy}
\end{equation}
This estimate also holds with any subset of the moduli on the left.

Write \(e_q(v)=\exp(2\pi i v/q)\).
In a separated Type II block, let \(d\) denote the variable of length
\(T\) whose reciprocal is placed in the moment and let \(f\) denote the
other variable, of length \(Y\). The modulus variable \(b\) is of size
\(X\); the uniform estimate below also applies with any other modulus
length in place of \(X\).
The block has the form
\[
 \sum_{b\asymp X}\beta_b'
 \sum_{\substack{d\asymp T,\ f\asymp Y,\ h\asymp J\\(df,b)=1}}
       a_d c_f w_h\, e_b(Nh\bar d\bar f),
\]
where all weights are bounded by \(O_\varepsilon(N^\varepsilon)\).
Signed frequencies can be split according to the sign of \(h\); it
suffices to consider one sign with \(\lvert h\rvert\asymp J\).
For each divisor \(g\mid N\), let \(S_g\) be the part with \((b,N)=g\).
Write \(b=gq\) and \(Q=X/g\). If a prime divides \(q\), its valuation in
\(b\) exceeds its valuation in \(g\), so its valuation in \(N\) equals
its valuation in \(g\). Hence \((N/g,q)=1\).
The character reduces as
\[
  e_b(Nh\bar d\bar f)=e_q((N/g)h\bar d\bar f)
\]
with the right-hand inverses taken modulo \(q\). For \(q=1\) the
character is identically one and all congruences have their single
residue class. The dyadic interval for \(q\) contains only positive
integers and is empty unless \(Q\gg1\).

For a fixed \(g,q\), extend the weights by zero to enforce
\((d,b)=(f,b)=1\), and define functions on \(\mathbb Z/q\mathbb Z\) by
\[
 A_q(z)=\sum_{\substack{d\asymp T\\(d,b)=1}}
                a_d e_q((N/g)z\bar d),\qquad
 r_q(z)=
  \sum_{\substack{\lvert h\rvert\asymp J,\ f\asymp Y\\
                   (f,b)=1,\ h\bar f\equiv z\;(\mathrm{mod}\ q)}} w_h c_f .
\]
Here the sum for \(r_q\) uses the chosen sign of \(h\).
The block for \(b\) equals \(\beta_b'\sum_z r_q(z)A_q(z)\).
With counting-measure norms, H\"older's inequality gives the bound
\[
 |\beta_b'|\|r_q\|_{2k/(2k-1)}\|A_q\|_{2k}.
\]
Expand \(\lvert A_q(z)\rvert^{2k}\) and sum over \(z\).
Orthogonality on \(\mathbb Z/q\mathbb Z\) imposes
\[
 (N/g)(\bar d_1+\cdots+\bar d_k-
 \bar d_{k+1}-\cdots-\bar d_{2k})\equiv0\pmod q.
\]
The coprimality of \(N/g\) to \(q\) removes this factor. Since each
\(d_i\) is coprime to \(b\) and hence to \(q\), taking absolute values
of the surviving weight products gives
\begin{equation}
 \|A_q\|_{2k}^{2k}\ll_\varepsilon N^\varepsilon q\,E_{q,k}(T).
 \label{eq:target-orthog-moment}
\end{equation}

The triangle inequality gives
\(\|r_q\|_1\ll_\varepsilon N^\varepsilon JY\).
Expansion of \(\|r_q\|_2^2\) requires
\(h_1\bar f_1\equiv h_2\bar f_2\pmod q\), equivalently
\(h_1 f_2\equiv h_2 f_1\pmod q\), since \(f_1,f_2\) are units.
The two integer products are nonzero and lie in an interval of length
\(O(JY)\). For each of the \(O(JY)\) possible first products there are
\(O(1+JY/q)\) possible second products in its residue class. A given
nonzero product \(m\) has \(O(\tau(|m|))\) ordered factorizations of
the specified sign. Bounding the weights and dropping the coprimality
conditions after the congruence has been converted into products yields
\[
 \|r_q\|_2^2\ll_\varepsilon N^\varepsilon JY(1+JY/q).
\]
To interpolate, put \(p_k=2k/(2k-1)\) and write
\[
 \lvert r_q\rvert^{p_k}=
 (\lvert r_q\rvert)^{p_k(1-1/k)}
 (\lvert r_q\rvert^2)^{p_k/(2k)}.
\]
The two exponents on the factors add to one. H\"older and the
\(p_k\)-th root then give
\[
 \|r_q\|_{p_k}\le \|r_q\|_1^{1-1/k}\|r_q\|_2^{1/k}
 \ll_\varepsilon N^\varepsilon
          (JY)^{1-1/(2k)}(1+JY/q)^{1/(2k)} .
\]
In particular the two interpolation inequalities are
\(\|r_q\|_{6/5}\le\|r_q\|_1^{2/3}\|r_q\|_2^{1/3}\) for \(k=3\)
and \(\|r_q\|_{4/3}\le\|r_q\|_1^{1/2}\|r_q\|_2^{1/2}\) for \(k=2\).

Combining with \eqref{eq:target-orthog-moment} gives a factor
\((q+JY)^{1/(2k)}E_{q,k}(T)^{1/(2k)}\) for each modulus,
in addition to \(N^\varepsilon(JY)^{1-1/(2k)}\).
H\"older in \(q\asymp Q\) gives
\[
 \sum_{q\asymp Q} E_{q,k}(T)^{1/(2k)}
 \ll Q^{1-1/(2k)}
         \left(\sum_{q\asymp Q}E_{q,k}(T)\right)^{1/(2k)} .
\]
Therefore \eqref{eq:target-mean-energy} yields
\begin{equation}
 \lvert S_g\rvert \ll_\varepsilon N^\varepsilon
    (JY)^{1-1/(2k)}(Q+JY)^{1/(2k)}
    Q^{1-1/(2k)}(QT^k+T^{2k})^{1/(2k)} .
 \label{eq:target-holder-moduli}
\end{equation}
For \(u,v\ge0\) and \(0<t\le1\), \((u+v)^t\le u^t+v^t\).
Applying this to both binomials in \eqref{eq:target-holder-moduli}
gives the following four terms, with \(N^\varepsilon\) suppressed.
\begin{equation}
 \begin{gathered}
 (JY)^{1-1/(2k)} Q^{1+1/(2k)}T^{1/2},\qquad
 (JY)^{1-1/(2k)} QT,\\
 JYQT^{1/2},\qquad
 JYQ^{1-1/(2k)}T .
 \end{gathered}
 \label{eq:target-moment-four-terms}
\end{equation}
The powers of \(Q\) are \(7/6,1,1,5/6\) in the sixth moment and
\(5/4,1,1,3/4\) in the fourth moment. They are all positive. Replacing
\(Q\) by \(X\) thus gives an upper bound, and the sum over \(g\mid N\)
adds at most \(\tau(N)\ll_\varepsilon N^\varepsilon\). In particular
the argument uses the same bounds when \(N\) is a primorial or has
arbitrarily high prime-power factors.

For the reversed sixth moment take \(T=L=N^s\) and
\(Y=W=N^{x-s}\). Then \(r_q\) groups \((h,n)\) by \(h\bar n\), and
the first factor in \eqref{eq:target-holder-moduli} is
\((JW)^{5/6}\). Upon dividing the four terms in
\eqref{eq:target-moment-four-terms} by \(X^2\), their powers of \(N\)
are, in order,
\begin{equation}
 \frac{5j-2s}{6},\qquad
 \frac{5j+s-x}{6},\qquad
 j-\frac{s}{2},\qquad
 j-\frac{x}{6}.
 \label{eq:target-sixth-rows}
\end{equation}
Indeed the unnormalized terms are
\((JW)^{5/6}X^{7/6}L^{1/2}\), \((JW)^{5/6}XL\),
\(JWX L^{1/2}\), and \(JWX^{5/6}L\).
For the fourth moment take \(T=W\), \(Y=L\) and group \((h,\ell)\) by
\(h\bar\ell\). The unnormalized terms are
\((JL)^{3/4}X^{5/4}W^{1/2}\), \((JL)^{3/4}XW\),
\(JLXW^{1/2}\), and \(JLX^{3/4}W\). Substitution of \(X=N^x\),
\(J=N^j\) and division by \(X^2\) now gives
\begin{equation}
 \frac{3j+s-x}{4},\qquad
 \frac{3j-s}{4},\qquad
 j+\frac{s}{2}-\frac{x}{2},\qquad
 j-\frac{x}{4}.
 \label{eq:target-fourth-rows}
\end{equation}

\section{Choice of cutoffs and the balanced thresholds}
\label{sec:exponents}

For \(N>1\) write \(X=N^x\), \(R=N^r\), \(L=N^s\), \(J=N^j\) and
\(H=N^{j_0}\). The case \(N=1\) is finite and obeys both bounds by
increasing the constants. For real cutoffs, the inequalities on
integer variables mean truncation at their integer parts. Here are
the exponent bounds relative to \(X^2\), each up to \(N^\varepsilon\).
The second-moment alternatives differ only in their nonzero
completed frequency.
\begin{equation}\label{eq:row-ledger}
\begin{array}{@{}c|l@{}}
 \text{Type I} & j-x,\quad j+2r-x/2\\
 \text{Elementary }2 & (j-s)/2,\quad j-x/2,\quad j+s/2-x/4\\
 \text{Spectral }2 & (j-s)/2,\quad j-x/2,\quad j+3s/8-x/4\\
 \text{Reversed }6 &
  (5j-2s)/6,\quad (5j+s-x)/6,\quad j-s/2,\quad j-x/6\\
 \text{Moment }4 &
  (3j+s-x)/4,\quad(3j-s)/4,\quad j+s/2-x/2,\quad j-x/4.
\end{array}
\end{equation}
The Type I and elementary second-moment rows come from the absolute
bounds of Section~\ref{sec:completion}. The spectral row comes from
Section~\ref{sec:spectral-completion} and the last two rows from
\eqref{eq:target-sixth-rows} and \eqref{eq:target-fourth-rows}.

Let \(0\le j\le c\), with \(c\) a cap. For \(d=9\) in the
elementary argument and \(d=8\) in the spectral argument assume
\(dc\le x\). Set
\[
 r=\frac{x-c}{5},\qquad u=\frac{2x-c}{5},\qquad
 t_9=\frac{3x-2c}{10},\qquad t_8=\frac{6x-4c}{17}.
\]
Use Type I and three Type II intervals, respectively
\([r,t_d]\), \([t_d,u]\), and \([u,x/2]\). Use the appropriate
second moment, the reversed sixth moment, and the fourth moment
on them. The gaps \(t_9-r=u-t_9=x/10\), 
\(t_8-r=(13x-3c)/85\), \(u-t_8=(4x+3c)/85\), and
\(x/2-u=(x+2c)/10\) are positive. Moreover
\(r-c=(x-6c)/5\ge x/20\) for both arguments. In particular
\(L\) in the spectral band exceeds the cap scale by a fixed power.

We show that every used exponent is at most \((dc-x)/10\).
All exponents increase in \(j\), so set \(j=c\).
On each interval every row is affine in \(s\).
Its maximum is attained at the upper endpoint for positive slope
in \(s\), the lower endpoint for negative slope, and either
endpoint for zero slope. The selected values are affine in \(x,c\).
Divide by \(x>0\) and test \(c/x=0,1/d\), since their differences
from the proposed bound are also affine. The entries below are
the selected values at \(x=1\), in the order of
\eqref{eq:row-ledger}.
\[
\begin{array}{c|c|c}
 d=9&c=0&c=1/9\\ \hline
 \mathrm{I}&-1,\ -1/10&-8/9,\ -1/30\\
 \mathrm{II},2&-1/10,\ -1/2,\ -1/10&-1/30,\ -7/18,\ 0\\
 \mathrm{II},6&-1/10,\ -1/10,\ -3/20,\ -1/6&
             0,\ -1/90,\ -1/36,\ -1/18\\
 \mathrm{II},4&-1/8,\ -1/10,\ -1/4,\ -1/4&
             -1/24,\ -1/90,\ -5/36,\ -5/36
\end{array}
\]
\[
\begin{array}{c|c|c}
 d=8&c=0&c=1/8\\ \hline
 \mathrm{I}&-1,\ -1/10&-7/8,\ -1/40\\
 \mathrm{II},2&-1/10,\ -1/2,\ -2/17&-1/40,\ -3/8,\ -1/272\\
 \mathrm{II},6&-2/17,\ -1/10,\ -3/17,\ -1/6&
             -1/272,\ 0,\ -5/136,\ -1/24\\
 \mathrm{II},4&-1/8,\ -1/10,\ -1/4,\ -1/4&
             -1/32,\ 0,\ -1/8,\ -1/8
\end{array}
\]
Each number is at most \(-1/10\) in the first column and at most
zero in the second. For example, the elementary transition rows at
\(t_9\) are \(c+t_9/2-x/4=(9c-x)/10=(5c-2t_9)/6\).
The spectral bottleneck is reached at \(u\), where
\((5c+u-x)/6=(3c-u)/4=(8c-x)/10\).
Fixed dyadic comparability in the intervals changes the row
bounds only by bounded factors because all the exponents are
affine functions of the logarithms of the lengths. The same
fixed factors can be included in the hypotheses \(J\ll H\) and
\(L\ll X^{1/2}\) of the completed estimates.

For the two main theorems take \(x=1-\theta\) and \(c=j_0=1-2\theta\).
The elementary cap condition is strict for \(17\theta>8\), and the
spectral cap condition is strict for \(15\theta>7\).
Summing the dyadic blocks gives the exponents
\((8-17\theta)/10\) and \((7-15\theta)/10\), respectively.
This proves Theorems~\ref{thm:elementary} and~\ref{thm:spectral}.

For the fourth-moment-only result take \(c\le x/10\), keep
\(r=(x-c)/5\), and use the elementary second moment up to
\(t=(x-c)/3\) followed by the fourth moment.
Here \(t-r=2(x-c)/15>0\) and \(x/2-t=(x+2c)/6>0\).
The maxima in the three groups at \(j=c\) are the entries of
\[
 \begin{gathered}
 \big(c-x,(6c-x)/10\big),\qquad
 \big((6c-x)/10,c-x/2,(10c-x)/12\big),\\
 \big((3c-x/2)/4,(10c-x)/12,c-x/4,c-x/4\big).
 \end{gathered}
\]
Subtracting the distinct entries from \((10c-x)/12\), in the
order of first occurrence, gives
\[
 \begin{gathered}
 (11x-2c)/12,\quad (28c+2x)/120,\quad (5x-2c)/12,\\
 0,\quad (c+x/2)/12,\quad (x-c)/6.
 \end{gathered}
\]
They are nonnegative for \(0\le c\le x/10\).
Substitution \(c=1-2\theta\) gives \((9-19\theta)/12\).

\begin{proof}[Proof of Proposition~\ref{prop:obstruction}]
Take \(j=j_0=x/8\), \(s=3j=3x/8\).
A saving in Type I requires \(r<(x/2-j)/2=3x/16\);
thus this scale cannot be absorbed by raising the cutoff.
It is a Type II scale since \(s<x/2\).
The second entries of the sixth- and fourth-moment rows are
\((5j+s-x)/6=0\) and \((3j-s)/4=0\).
The spectral second-moment row gives
\(j+3s/8-x/4=x/64>0\). Hence none of the three available
Type II estimates has all its exponents negative. Any larger
frequency cap still includes \(j=x/8\), so this obstruction
precludes a uniform strict saving from the same estimates
there as well. The strict cap range was established above.
\end{proof}

\section{The unbalanced range of the elementary nine estimates}
\label{sec:unbalanced}

\begin{proof}[Proof of Proposition~\ref{prop:unbalanced}]
Let \(A=N^a\) be the length of the variable to which the M\"obius
identity is applied and let \(B=N^b\) be the modulus length. The
frequency cap is \(AB/N\), so its exponent is \(j_0=a+b-1\).
The normalized sum is measured in units of \(AB\).
For rectangles the unordered pair \(\{a,b\}\) is
\(\{1-p,1-q\}\). The orientation \(a=1-p\ge b=1-q\) is available
when the longer coefficient carries the M\"obius weight. In the
rectangle application both coefficients carry this weight.
The other orientation is obtained by the identity
\[
 e(Nh\bar m/n)=e(-Nh\bar n/m)\,e(Nh/(mn)),\qquad (m,n)=1.
\]
Indeed \(m\bar m+n\bar n-1\) is divisible by both \(m\) and \(n\),
hence by \(mn\). On the dyadic block, the extra factor \(e(Nh/(mn))\)
is a smooth weight with bounded derivatives in the scaled variables
when \(NJ\ll AB\), so it can be included in the Fourier separation.

We first record the nine quantities governing this argument. Write
\(R=N^r\), \(L=N^s\), \(W=A/L\), and \(J=N^j\).
The cutoff exponent satisfies \(r\le a/2\), and the Type II blocks
satisfy \(r\le s\le a/2\).
The powers of \(N\), after dividing the original sum by \(AB\), are
\begin{equation}
\begin{array}{@{}c|l@{}}
 \mathrm{I}
 & j-b,\quad j+2r+b/2-a\\[2pt]
 \mathrm{II},2
 & (j-s)/2,\quad j-b/2,\quad j+s/2+b/4-a/2\\[2pt]
 \mathrm{II},4
 & (3j+s+b-2a)/4,\quad (3j-s)/4,\quad
          j+s/2-a/2,\quad j-b/4 .
\end{array}
\label{eq:target-unbalanced-rows}
\end{equation}
Here and below fixed dyadic constants only change multiplicative
constants in the estimates.

To obtain the first line, use the completed length \(Y=A/d\) at a
Type I divisor \(d\le R^2\). The Ramanujan average contributes
\(N^\varepsilon JY\) after summing moduli for fixed \(d\). Summing
\(Y=A/d\) gives \(N^\varepsilon JA\). The bound for the nonzero frequencies
is \(N^\varepsilon B^{1/2}\) per modulus for each \(d,h\): the completion
factor \(Y/B\) cancels the Fourier length \(B/Y\). Summing over all
\(b\)-variables, \(h\)-variables and \(d\)-variables gives
\(N^\varepsilon JR^2B^{3/2}\).
For the second line, the outer Cauchy--Schwarz factor is
\(N^\varepsilon BW\). Before this factor is included, the three squared
contributions are bounded by \(N^\varepsilon BWJL\),
\(N^\varepsilon WJ^2L^2\), and \(N^\varepsilon B^{3/2}J^2L^2\),
respectively. The first is the rational diagonal, with
\(N^\varepsilon JL\) tuples and \(O(BW)\) summands for each tuple.
The second is the zero Poisson frequency away from the diagonal,
where the Ramanujan average over the moduli cancels the denominator
in \(W/B\). The third is the nonzero completion, for which the
length \(B/W\) cancels \(W/B\) and Weil contributes \(B^{1/2}\) per
modulus. Multiplying by \(BW\), taking square roots, and dividing by
\(AB\) gives the second line of \eqref{eq:target-unbalanced-rows}.
For the last line, use \eqref{eq:target-moment-four-terms} with
\(k=2\), \(T=W\), \(Y=L\), and replace \(Q\) by \(B\). The four
unnormalized terms are
\[
 (JL)^{3/4}B^{5/4}W^{1/2},\quad (JL)^{3/4}BW,\quad
 JLBW^{1/2},\quad JLB^{3/4}W.
\]
Since \(W=N^{a-s}\), division by \(AB\) gives the asserted powers.
Setting \(a=b=x\) in \eqref{eq:target-unbalanced-rows} recovers
the corresponding rows of \eqref{eq:row-ledger}.

The nonzero Type I row is unnecessary if \(a-2r>b\), because
\(Y/B\gg N^{a-2r-b}\) uniformly in \(d\le R^2\).
The nonzero Type II Cauchy--Schwarz row is unnecessary at \(s\) if
\(a-s>b\), because \(W/B\gg N^{a-s-b}\).
Indeed for a smooth compactly supported completion weight \(V\),
\(\widehat V(t)\ll_M(1+|t|)^{-M}\). If the completed scale divided
by \(B\) exceeds \(N^\delta\), then the sum over nonzero integer
frequencies is bounded by \(O_M(N^{-\delta M}\sum_{k\ne0}|k|^{-M})\).
The remaining factors are polynomial in \(N\), so any required
power saving follows by choosing \(M\). We use this observation with
strict positive exponent separation.

Each row in \eqref{eq:target-unbalanced-rows} increases with \(j\),
and the long-completion conditions just stated are independent of it.
We can therefore test the cap \(j=j_0\), and use the resulting cutoffs
at every smaller nonnegative \(j\). For the rest of the proof write
\(j=j_0\). Define
\[
 r_{\mathrm I}=\frac{a-b/2-j}{2},\qquad c=a-b/2-2j .
\]
In the presence of its nonzero row, Type I needs \(r<r_{\mathrm I}\).
The Cauchy--Schwarz rows alone give savings exactly when \(s>j\),
\(j<b/2\), and one of
\[
                  s<c,\qquad s<a-b
\]
holds. The first inequality controls the nonzero row and the second
gives rapid completion decay. At equality of lengths the rapid-decay
argument supplies no power, so the inequalities used here are strict.

Consider first \(a\ge b\). If the top \(s=a/2\) is covered by
Cauchy--Schwarz with the nonzero row, then necessarily
\(b>2j\) and \(a-b>4j\). These two conditions suffice to cover the entire
Type II range by Cauchy--Schwarz and to meet Type I. In fact
\[
 c>a/2,\qquad
 r_{\mathrm I}-j
   =\tfrac12(a-b/2-3j)>\tfrac14 a-\tfrac12j>0,
\]
where \(a-b>4j\) gives \(j<a/4\).
Choose \(j<r<\min(r_{\mathrm I},a/2)\).
The Type I zero row is negative since \(b>2j\), its nonzero row
is negative by the upper bound on \(r\), and all the Cauchy--Schwarz
rows are negative on \([r,a/2]\).
In rectangle coordinates the two conditions are
\[
                      2p+q>1,\qquad 3p+5q>4 .
\]

If the top is covered by rapid Cauchy--Schwarz completion, then
necessarily \(b>2j\) and \(a>2b\). These conditions give \(a/2>j\)
and \((a-b)/2>j\), since \(a-b>b>2j\).
Choose \(j<r<\min((a-b)/2,a/2)\). Every Type I completed length
exceeds \(B\) by a fixed power, and so does every Type II completed
length since \(a-a/2>b\). The Type I zero row and the two remaining
Cauchy--Schwarz rows are negative. These conditions become
\[
                        2p+q>1,\qquad 2q-p>1 .
\]

Suppose the top \(s=a/2\) is covered by the fourth moment.
Its last row requires \(b>4j\) and its second row requires
\(a/2>3j\). In particular \(c>a-b\), so the range \(s<a-b\) supplies
no further Cauchy--Schwarz blocks beyond those with \(s<c\).
At \(s=3j\) the fourth-moment second row is zero. If \(r\le3j\), a
uniform power saving here must come from Cauchy--Schwarz and
requires \(3j<c\). If \(r>3j\), then Type I requires \(r<r_{\mathrm I}\):
indeed the alternative \(r<(a-b)/2\) also implies \(r<r_{\mathrm I}\)
because \(b/2>j\). In this case
\(a-b/2>2r+j>7j\ge 5j\).
Both cases give the necessary pair
\[
                            b>4j,\qquad a-b/2>5j .
\]
Conversely, this pair gives \(a>7j\), since \(b/2>2j\).
For \(j=0\) it gives \(a>0\) as well. Thus
\(3j<a/2\), \(r_{\mathrm I}>j\), and \(c>3j\).
Choose \(t\) with \(3j<t<\min(c,a/2)\) and \(r\) with
\(j<r<\min(r_{\mathrm I},t)\). Type I is valid for this cutoff.
Cauchy--Schwarz gives savings on \([r,t]\).
For \(s\in[t,a/2]\), the second and fourth fourth-moment rows
are negative by \(t>3j\) and \(b>4j\).
The third is at most \(j-a/4<0\). The first is at most
\((3j+b-3a/2)/4\le (3j-a/2)/4<0\), using \(b\le a\).
These are exactly the conditions
\[
                       4p+3q>3,\qquad 8p+11q>9 .
\]
The three top cases exhaust the method. Raising \(r\) to or above
\(a/2\) cannot give a Type I saving on its own: the nonzero row is then
at least \(j+b/2\), and the long-completion condition is impossible.
Choosing strict cutoffs gives a uniformly negative maximum of the
relevant rows on the closed intervals. Conversely, a row tending to
zero as \(s\) approaches a transition supplies no fixed power saving
there, which accounts for the strict boundaries in the three cases.

For the reverse orientation, write \(a_{\mathrm r}\le b_{\mathrm r}\)
for the factorized and modulus exponents, respectively. The top cannot
be covered by Cauchy--Schwarz through either
\(a_{\mathrm r}-b_{\mathrm r}>4j\) or
\(a_{\mathrm r}>2b_{\mathrm r}\).
If it is covered by the fourth moment, the preceding necessity
argument, which did not assume an ordering of \(a,b\) in this case,
gives
\[
 b_{\mathrm r}>4j,\qquad a_{\mathrm r}-b_{\mathrm r}/2>5j .
\]
In particular \(a_{\mathrm r}>4j\) and
\(b_{\mathrm r}-a_{\mathrm r}/2\ge a_{\mathrm r}-b_{\mathrm r}/2>5j\).
Thus the original orientation is already in the mixed region. This
proves the exact union for the given estimates and the two choices of
orientation.

To make the comparison, apply Bettin--Chandee
\cite[Theorem~1]{BC} to lengths \(A,B,J\) with phase multiplier \(N\).
For coefficients of size \(N^\varepsilon\) the product of the three
\(\ell^2\) norms is \(O_\varepsilon(N^\varepsilon(ABJ)^{1/2})\).
At \(J\le AB/N\) the factor \((1+NJ/(AB))^{1/2}\) is bounded.
The two remaining terms are
\((ABJ)^{7/20+\varepsilon}(A+B)^{1/4}\) and
\((ABJ)^{3/8+\varepsilon}(JA+JB)^{1/8}\).
At \(a\ge b\), and after multiplication by the norms and division
by \(AB\), the two relative exponents are
\[
  (17j+2a-3b)/20,\qquad j-b/8 .
\]
Inserting \(j=1-p-q\), \(a=1-p\), \(b=1-q\) gives
\((16-19p-14q)/20\) and \((7-8p-7q)/8\). Both are negative exactly
in the stated Bettin--Chandee region.
The identities
\[
\begin{aligned}
 4p+3q&=\frac7{33}(19p+14q)+\frac{q-p}{33},\\
 8p+11q&=\frac{19}{33}(19p+14q)+\frac{97}{33}(q-p)
\end{aligned}
\]
show that \(19p+14q>16\) and \(q\ge p\) imply both inequalities of
the mixed region, since \(112/33>3\) and \(304/33>9\).
For strictness take \(p=q=12/25\). Then \(p+q<1\),
\(4p+3q=84/25>3\), and \(8p+11q=228/25>9\), while
\(19p+14q=396/25<16\).
\end{proof}

\section{The Goldbach remainder rectangles}\label{sec:rectangles}

We prove Corollary~\ref{cor:rectangles}. Write \(X=N^{1-\theta}\) and
\(H=X^2/N\). The implicit constants in \(U\asymp V\asymp N^\theta\)
are fixed. In the smooth rectangle the summand is multiplied by
\(W_1(u/U)W_2(v/V)\), with fixed compactly supported smooth functions
on the positive real axis. In particular the proof applies to smooth
weights supported in the dyadic rectangle. The sharp rectangle uses the
two indicators of \([1,2)\). All variables in the rectangle sum are
positive integers.

\subsection{Smoothing, truncation, and Poisson summation}

Let \(\rho,\sigma,\kappa>0\) be small constants to be chosen below.
In the sharp case put \(\xi=N^{-\sigma}\) and convolve each indicator
with \(\xi^{-1}\varphi(\,\cdot\,/\xi)\), where \(\varphi\) is a fixed
nonnegative smooth function of integral one supported in \((-1,1)\).
Each resulting weight is bounded by one, is supported in
\([1-\xi,2+\xi]\), and agrees with the indicator outside the two
intervals of radius \(\xi\) about its endpoints. Its \(i\)-th
derivative is \(O_i(\xi^{-i})\). For a fixed \(u\asymp U\) the absolute
number of solutions is at most
\[
 \sum_{1\le a<N/u}\tau(N-au)\ll_\varepsilon N^\varepsilon N/U.
\]
There are \(O(\xi U+1)\) such integers \(u\) in the two transition
intervals. Fixing \(v\) gives the same bound with \(V\) in place of
\(U\). Applying these two estimates to the difference of the products
of the weights gives the error
\begin{equation}\label{eq:rectangle-smoothing}
 O_\varepsilon\bigl(N^{1+\varepsilon}(\xi+U^{-1}+V^{-1})\bigr).
\end{equation}
The derivative bounds for these approximants also hold for the fixed
smooth weights, with constants depending on them.

Put \(d_0=(a,b)\). Then \(d_0\mid N\). For each \(d_0\ge N^\kappa\)
we can count solutions by dropping the coprimality condition on
\(a/d_0,b/d_0\) and all interval restrictions. With \(P=N/d_0\) the
resulting count is at most
\[
 \sum_{m=1}^{P-1}\tau(m)\tau(P-m)\ll_\varepsilon P N^\varepsilon.
\]
The sum over these divisors is
\(O_\varepsilon(N^{1-\kappa+\varepsilon})\), with the divisor factor
absorbed by relabeling \(\varepsilon\). Choose a fixed smooth
\(0\le\chi\le1\) that is zero on \((-\infty,1]\) and one on
\([2,\infty)\). Multiply the remaining summands by
\(\chi(au/N^{1-\rho})\chi(bv/N^{1-\rho})\). For
\(m=au<2N^{1-\rho}\) the number discarded is bounded by
\(\sum_m \tau(m)\tau(N-m)\ll_\varepsilon N^{1-\rho+\varepsilon}\).
For \(m=bv<2N^{1-\rho}\) the same displayed divisor sum counts the
discarded terms with the roles of the two products interchanged.
These estimates apply to the smooth weights as well, after
multiplication by their sup norms.

For a fixed remaining \(d_0\mid N\), write
\[
 a=d_0 a',\quad b=d_0b',\quad P=N/d_0,\quad Z=X/d_0,\quad
 (a',b')=1.
\]
The support conditions give \(a',b'\ll Z\) and
\(a',b'\gg ZN^{-\rho}\), since \(u\asymp U\), \(v\asymp V\) and
\(au,bv\ge N^{1-\rho}\) on that support. Insert a smooth dyadic
partition in \(a'\) and \(b'\). It has \(O((\log N)^2)\) pairs of
nonempty pieces, with parameters \(A,B\) in
\([cZN^{-\rho},C Z]\) and cutoff factors
\(\omega_A(a'/A),\omega_B(b'/B)\) supported in a fixed compact
subset of \((0,\infty)\), with uniformly bounded derivatives.
For example a dyadic partition on positive reals is obtained by
summing the dilates of \(\varpi(t)-\varpi(2t)\), where \(\varpi=1\)
on \([0,1]\) and \(\varpi=0\) on \([2,\infty)\).

For real \(t\), define
\[
 \begin{split}
 G(t;x_1,x_2)={}&W_1(t)\,
 W_2\!\left(\frac{N-d_0 A x_1Ut}{d_0 Bx_2 V}\right)\\
 &{}\times
 \chi\!\left(\frac{d_0 A x_1Ut}{N^{1-\rho}}\right)
 \chi\!\left(\frac{N-d_0 A x_1Ut}{N^{1-\rho}}\right),
 \end{split}
\]
and use the Fourier convention
\(\widehat G(z;x_1,x_2)=\int_{\mathbb R}G(t;x_1,x_2)e(-zt)\,dt\).
The equation \(a'u+b'v=P\) gives \(u\equiv P\overline{a'}\pmod{b'}\)
and determines \(v\) from \(u\). The smooth weights enforce positivity.
Poisson summation for this residue class yields the exact \(u\)-sum
\begin{equation}\label{eq:rectangle-outer-poisson}
 \frac U{b'}\sum_{h\in\mathbb Z}
 e\!\left(\frac{hP\overline{a'}}{b'}\right)
 \widehat G\!\left(\frac{hU}{b'};\frac{a'}A,\frac{b'}B\right).
\end{equation}
Indeed the residue-class indicator is the average of additive
characters of modulus \(b'\), and Fourier transformation of
\(G(u/U;x_1,x_2)\) contributes \(U\widehat G(U\,\cdot\,;x_1,x_2)\).
This gives both the factor \(U/b'\) and the plus sign in the phase of
\eqref{eq:rectangle-outer-poisson}.

For \(x_1,x_2\) in the cutoffs and \(t\) in the fixed support of \(W_1\),
all derivatives of each inner argument in \(G\) of any specified
order are \(O_i(N^\rho)\). For the argument of \(W_2\) this follows
from \(N/(d_0 B V)\ll N^\rho\), \(d_0 A U\ll N\), and \(x_2\)
bounded away from zero. For the last two arguments it follows from
\(d_0 A U\ll N\). The chain and product rules give derivative
bounds of total order \(i\) by \(O_i(N^{C_0i(\rho+\sigma)})\)
for an absolute \(C_0\). Integration by parts in \(t\) gives, for
each integer \(m\ge2\),
\begin{equation}\label{eq:rectangle-poisson-decay}
  |\widehat G(z;x_1,x_2)|
  \ll_m N^{C_0m(\rho+\sigma)}(1+|z|)^{-m}.
\end{equation}
Choose an auxiliary smooth even frequency cutoff equal to one
for \(|h|U/(BN^\delta)\le1\) and zero for that ratio at least two,
where \(\delta>C_0(\rho+\sigma)\). The part removed is negligible.
To see uniformity even when the cutoff contains no nonzero integer,
put \(v_0=U/B\). For each sign, the tail uses
\(|h|v_0>N^\delta\) and
\[
 \sum_{\substack{h\ge1\\hv_0>N^\delta}}(1+hv_0)^{-m}
 \ll_m (1+v_0^{-1})N^{-(m-1)\delta}.
\]
The inequality follows by bounding the first term and comparing the
remaining decreasing sequence with its integral. Multiplication by
\(U/b'\asymp v_0\) costs \(O(1+v_0)\), a fixed polynomial in \(N\).
There are at most a fixed power of \(N\) choices of \(d_0,a',b'\)
and the dyadic pieces. Taking \(m\) sufficiently large gives an error
\(O_M(N^{-M})\) for any fixed \(M>0\).
If \(d_0>H\), then \(B/U\ll H/d_0< O(1)\). The same tail estimate
still applies, and there are no retained nonzero terms when
\(d_0\) exceeds a constant multiple of \(HN^\delta\).
Divisors between \(H\) and this inflated upper limit are included
in the following argument.

\subsection{The nonzero frequencies}

Insert a smooth dyadic partition in \(|h|\ge1\). For either sign and
each nonempty piece with \(|h|\asymp J\), \(J\ge1\), the sum has
prefactor \(U/B\) and a smooth weight in the normalized coordinates
\((x_1,x_2,z)=(a'/A,b'/B,|h|/J)\). This weight is explicitly the
product of the three dyadic cutoffs, \(x_2^{-1}\), the frequency
cutoff at \(JzU/(BN^\delta)\), and
\(\widehat G(\pm JzU/(Bx_2);x_1,x_2)\).
Every derivative of bounded order \(i\) is
\(O_i(N^{C_i(\rho+\sigma+\delta)})\). In fact a derivative on the
exponential in the defining integral for \(\widehat G\) costs at
most \(O_i(1+JU/B)\ll N^\delta\), while derivatives on \(G\)
obey the bound preceding \eqref{eq:rectangle-poisson-decay}.
Derivatives on the two frequency cutoffs are bounded because
\(z\) is in a fixed compact subset of \((0,\infty)\) and
\(JU/(BN^\delta)\ll1\) for a surviving piece. Fourier separation
of a smooth compactly supported weight uses a fixed finite number of
these derivative bounds. The two signs give the same estimates, by
replacing the phase by its conjugate and conjugating the weights.

For nonsquarefree \(d_0\) the M\"obius product is zero. For squarefree
\(d_0\),
\begin{equation}\label{eq:rectangle-moebius-gcd}
 \mu(d_0a')\mu(d_0b')
 =\mu(a')\mu(b')\mathbf 1_{(a',d_0)=1}\mathbf 1_{(b',d_0)=1}.
\end{equation}
The \(b'\)-factor is an admissible bounded coefficient. Here is the
effect of the \(a'\)-restriction on the row estimates
\eqref{eq:row-ledger}. In Type~I, after the M\"obius split, fix the
short convolution variable \(\nu\) and write \(a'=\nu m\) with
\(m\asymp Y=A/\nu\). Use
\[
 \mathbf 1_{(a',d_0)=1}
 =\sum_{\substack{\lambda\mid d_0\\\lambda\mid a'}}\mu(\lambda).
\]
For a fixed \(\lambda\), put \(\lambda'=\lambda/(\lambda,\nu)\).
The divisibility condition is precisely \(m=\lambda' m_1\).
The phase is then \(e(hP\overline{\nu\lambda'}\overline{m_1}/b')\)
on the unit class, and the \(m_1\)-sum is at scale
\(Y'=Y/\lambda'\). Moduli with \((b',\nu\lambda')>1\) contribute
no terms. The zero Poisson frequency has factor \(Y'/b'\le Y/b'\)
and the Ramanujan bound with numerator \(hP\), since
\(\nu\lambda'\) is a unit. In the nonzero frequencies the Weil
estimate is used with \(K=B/Y'\) and the length cancels by
\(Y'K=B\). Thus shortening \(Y\) has no effect on the resulting
bound. There are at most \(\tau(d_0)\) terms in this inclusion-exclusion
sum. In Type~II, \(a'=\ell n\) satisfies \((a',d_0)=1\) if and
only if both \((\ell,d_0)=1\) and \((n,d_0)=1\). Multiplying the
two separated coefficient sequences by these indicators preserves
their upper bounds in Cauchy--Schwarz and the two moment estimates.

For clarity, the scale parameters in the row estimates can be
bounded using \(Z\). The Type~I absolute bounds are the sum of
\(O_\varepsilon(JA N^\varepsilon)\) and
\(O_\varepsilon(JR^2 B^{3/2}N^\varepsilon)\), which increase
when \(A\) and \(B\) are replaced by their upper scale \(Z\).
The support here is above \(R\), so the short term of the M\"obius
identity vanishes. In elementary Type~II write \(A\asymp LW\)
and take \(L\ll \sqrt Z\). The diagonal, Ramanujan and nonzero
Weil contributions to the completed \((b',n)\) second moment
are bounded, respectively, by constants times
\[
   BWJL,\qquad WJ^2L^2 N^\varepsilon,\qquad
   B^{3/2} J^2 L^2 N^\varepsilon.
\]
The Cauchy--Schwarz factor squared is \(BW N^\varepsilon\).
All these expressions increase under \(B\le C Z\) and
\(W\ll Z/L\). In the fourth and sixth moments the outer bounds
use nonnegative powers of both \(L,W\) and strictly positive
powers of the reduced modulus length \(Q\), so \(W\ll Z/L\)
and \(Q\ll Z\) give the indicated balanced rows at \(Z\).
The spectral row uses the rectangular estimate
\eqref{eq:spectral-rectangular} with \(A,B\gg ZN^{-\rho}\) and loses only
\(N^{O(\rho+\delta)+\varepsilon}\). All analytic estimates here
apply to \(P=N/d_0\) with divisor and Fourier losses measured
against \(N\).

Set \(x=1-\theta\), \(j_0=1-2\theta\) and \(w=\log d_0/\log N\).
The power corresponding to \(Z\) is \(x-w\).
The truncated nonzero frequencies, including the dyadic factors,
satisfy
\[
  J\ll (B/U)N^\delta\ll N^{j_0-w+\delta}.
\]
For large \(N\) we can use the cap \(c=j_0-w+2\delta\).
On every surviving block \(c\ge0\), and \(Z^2/P=H/d_0\le N^c\).
Put \(k_*=9\) in the elementary argument and \(k_*=8\) in the
spectral argument. For the asserted range with \(\theta\le1/2\)
one has \(\gamma=(x-k_*j_0)/10>0\). Upon taking
\(2k_*\delta<10\gamma\), the exact identities
\[
 \begin{split}
 (x-w)-k_*c&=10\gamma+(k_*-1)w-2k_*\delta>0,\\
 \frac{k_*c-(x-w)}{10}&=-\gamma+\frac{2k_*\delta-(k_*-1)w}{10}
 \end{split}
\]
give the cap condition and exponent from
Section~\ref{sec:exponents}. The gain in the numerator of the
exponent is \(8w\) in the elementary case and \(7w\) in the spectral
case. Constant factors in the dyadic ranges change log-base-\(N\)
exponents by \(O(1/\log N)\), or can be covered by an arbitrarily
small fixed part of the final \(\varepsilon\). For a fixed
\(d_0,A,B,J\) the nonzero contribution is therefore
\[
 \ll_\varepsilon \frac{U}{B}Z^2
 N^{-\gamma+(2k_*\delta-(k_*-1)w)/10
       +O(\rho+\sigma+\delta)+\varepsilon}.
\]
Since \(UZ/B\ll U N^\rho\), the prefactor is
\(O(N^{1-w+\rho})\). The sum over all pieces and divisors thus
saves a fixed power of \(N\) after the small parameters are chosen.

\subsection{The zero frequency under RH}

Let \(\mu_q(m)=\mu(m)\mathbf 1_{(m,q)=1}\). Under RH,
\begin{equation}\label{eq:rectangle-coprime-mertens}
  \sum_{m\le t}\mu_q(m)
  \ll_\varepsilon t^{1/2+\varepsilon}q^\varepsilon
  \qquad(t\ge1,\ q\ge1).
\end{equation}
Here is the uniformity in \(q\). The untwisted bound
\(\sum_{m\le t}\mu(m)\ll_\varepsilon t^{1/2+\varepsilon}\)
follows from Littlewood's convergence of
\(\sum_m\mu(m)m^{-s}\) for \(\operatorname{Re}s>1/2\) under
RH~\cite{littlewood1912}. If
\(F(t)=\sum_{m\le t}\mu(m)m^{-(1/2+\varepsilon)}\), this convergence
makes \(F\) bounded and partial summation gives the desired untwisted
bound. Let \(1_{q^\infty}(m)\) be the indicator that every prime
dividing \(m\) divides \(q\). Then \(\mu_q=\mu*1_{q^\infty}\)
by comparison of the local factors at each prime. Consequently
\[
 \left|\sum_{m\le t}\mu_q(m)\right|
 \ll_\varepsilon t^{1/2+\varepsilon}
      \prod_{p\mid q}(1-p^{-1/2-\varepsilon})^{-1}.
\]
For primes above a constant depending only on \(\varepsilon\) the
factor at \(p\) is at most \(p^\varepsilon\). The remaining factors
have a bounded product. This proves \eqref{eq:rectangle-coprime-mertens}.

For fixed \(b'\asymp B\) use \(q=b'd_0\) in
\eqref{eq:rectangle-coprime-mertens}. In the \(h=0\) term of
\eqref{eq:rectangle-outer-poisson}, the \(a'\)-weight at scale
\(A\) is \(\omega_A(a'/A)\widehat G(0;a'/A,b'/B)\).
Its sup norm and its total variation are
\(O(N^{C(\rho+\sigma)})\), by the product and derivative estimates
above. Partial summation on \(a'\) and \(b'd_0\ll N\) give
an upper bound \(O_\varepsilon(A^{1/2}N^{C(\rho+\sigma)+\varepsilon})\)
for that weighted sum. The factor \(U/b'\) and the \(O(B)\) choices
of \(b'\) give
\[
 \ll_\varepsilon U A^{1/2} N^{C(\rho+\sigma)+\varepsilon}
 \ll_\varepsilon
 N^{(1+\theta)/2+C(\rho+\sigma)+\varepsilon}.
\]
This also bounds the sum of the zero modes over \(d_0,A,B\), with
an arbitrarily small further \(N^\varepsilon\) loss. It uses RH
for \(\zeta\) alone.

Choose \(\kappa>0\) smaller than \(x/4\).
Fix \(\varepsilon>0\) sufficiently small compared with \(\gamma\)
and \(1-\theta\) when \(\theta\le1/2\).
After this choice all required derivative orders in the row
estimates are fixed. Choose \(\delta>0\) so small that the
loss \(O(\delta)+2k_*\delta/10\) is less than \(\gamma/4\).
Then take \(\rho,\sigma>0\) small enough that
\(C_0(\rho+\sigma)<\delta\), that the derivative losses leave
strict margins in both the zero and nonzero estimates, and that
\(\rho<x/4\). In the preliminary errors the divisor exponent can be
chosen independently of the analytic \(\varepsilon\). Take it below
\(\min(\kappa,\rho,\sigma,\theta)/2\).
All terms are now \(O(N^{1-\eta})\) for some \(\eta>0\) depending on
\(\theta\), including \eqref{eq:rectangle-smoothing} and the omitted
gcd and endpoint terms. The condition \(\rho<x/4\) also ensures
\(A\gg R\): for \(d_0<N^\kappa\) the cutoff exponent is
\((x-w-c)/5\le (x-w)/5\), whereas
\(A\gg N^{x-w-\rho}\).

For \(1/2<\theta<1\) the same argument requires only the zero mode.
Choose \(0<\delta<2\theta-1\) and
\(C_0(\rho+\sigma)<\delta\). Then \(BN^\delta/U\ll
N^{1-2\theta+\delta}\) tends to zero for every \(d_0\) and
every dyadic box. The nonzero modes all lie in the
rapidly decaying tail, and the zero bound still gives a strict
saving because \((1+\theta)/2<1\). The small parameters can be
chosen with \(C(\rho+\sigma)+\varepsilon<(1-\theta)/2\)
and with positive margins in all preliminary errors.
\qed

\section{Formal Verification Scope}
\label{sec:formal}
\begingroup
\raggedright

The accompanying \texttt{route-f} development uses Lean~4.36.0 and
Mathlib~v4.36.0. The modules below live under \texttt{lean/FareyGoldbach/}.
The repository's audit file \path{lean/tests/RouteFAudit.lean} uses
\texttt{\#print axioms} for the named declarations. This audit passes with
warnings treated as errors. The 38 named Target~3B declarations have no
\texttt{sorry} dependencies and use only the standard axioms
\texttt{propext}, \path{Classical.choice}, and \path{Quot.sound}.

\subsection*{\texttt{Target3BMoebiusSplit}}
The convolution and finite split are
\path{Target3B.moebius_two_cutoffs} and
\path{Target3B.moebius_finite_split}. The support assertions are
\path{Target3B.typeI_vanish_of_mul_cutoffs_lt} and
\path{Target3B.typeII_tail_vanish_of_le_cutoff}.
For the fourth-moment-only argument, the audited exponent
declarations include
\path{Target3B.typeI_scale},
\path{Target3B.typeII_completion_scales},
\path{Target3B.typeII_fourth_scales},
\path{Target3B.every_row_strict},
\path{Target3B.sum_rows_strict},
\path{Target3B.transition_optimal},
\path{Target3B.truncation_optimal}, and
\path{Target3B.threshold}. The related
\path{Target3B.analytic_power_bound} takes
\path{Target3B.AnalyticEstimates} as a hypothesis. Its branch
inequalities are analytic input to that implication.

\subsection*{\texttt{Target3BKaratsuba}}
The signed unit-fraction divisibility and its size consequence are
\path{Target3B.karatsuba_lcm_sq_dvd_prod} and
\path{Target3B.karatsuba_lcm_sq_le_prod}. The specializations are
\path{Target3B.karatsuba_four_reciprocals} and
\path{Target3B.karatsuba_six_reciprocals}. These are unconditional
kernel-checked statements.

\subsection*{\texttt{Target3BThresholds}}
In namespace \path{ThresholdEightSeventeenths}, the audited
declarations are
\path{cuts_ordered}, \path{type_i_rows}, \path{cauchy_rows},
\path{sixth_moment_rows}, \path{fourth_moment_rows},
\path{attained_at_sixth_cut}, \path{balanced_threshold}, and
\path{shifted_poisson_cap}.
In namespace \path{ThresholdSevenFifteenths}, they are
\path{cuts_ordered}, \path{type_i_rows}, \path{spectral_cauchy_rows},
\path{spectral_length_margin}, \path{exceptional_factor_exponents},
\path{normalize_completed_terms}, \path{sixth_moment_rows},
\path{fourth_moment_rows}, \path{attained_at_fourth_cut},
\path{obstruction_at_three_j}, \path{spectral_strictly_improves},
\path{balanced_threshold}, and \path{shifted_poisson_cap}.
These declarations check the rational exponent closures, the
spectral length and normalization inequalities, the change of cap
under division by \(d_0\), and the exact obstruction. The executable
\path{notes/check_target3b_thresholds.py} checks the corresponding
affine vertex inequalities using exact rational arithmetic.

The analytic derivations in the preceding sections are paper proofs.
Their external inputs are the full-modulus Weil bound, Pascadi's
local spectral estimate, and, for the outer zero frequency of the
rectangle application, the RH bound for the M\"obius summatory
function. The Lean statements cover the algebraic and rational
steps described above. The numerical feasible region for unbalanced
rectangles is proved in Section~\ref{sec:unbalanced}.
\par\endgroup

\begin{thebibliography}{10}
\raggedright

\bibitem{BC}
S.~Bettin and V.~Chandee,
\emph{Trilinear forms with Kloosterman fractions},
Adv. Math. \textbf{328} (2018), 1234--1262.
\href{https://doi.org/10.1016/j.aim.2018.01.026}{doi:10.1016/j.aim.2018.01.026}.

\bibitem{BGrings}
J.~Bourgain and M.~Z.~Garaev,
\emph{Kloosterman sums in residue rings},
Acta Arith. \textbf{164} (2014), 43--64.
\href{https://doi.org/10.4064/aa164-1-4}{doi:10.4064/aa164-1-4}.

\bibitem{BGprime}
J.~Bourgain and M.~Z.~Garaev,
\emph{Sumsets of reciprocals in prime fields and multilinear Kloosterman sums},
Izv. Math. \textbf{78} (2014), 656--707.
\href{https://doi.org/10.1070/IM2014v078n04ABEH002703}{doi:10.1070/IM2014v078n04ABEH002703}.

\bibitem{DRZ}
A.~Dong, N.~Robles, and D.~Zeindler,
\emph{Bilinear forms with Kloosterman fractions and applications},
preprint, \href{https://arxiv.org/abs/2601.00292v3}{arXiv:2601.00292v3}, 2026.

\bibitem{FKM}
\'E.~Fouvry, E.~Kowalski, and P.~Michel,
\emph{Algebraic trace functions over the primes},
Duke Math. J. \textbf{163} (2014), 1683--1736.
\href{https://arxiv.org/abs/1211.6043}{arXiv:1211.6043}.

\bibitem{IK}
H.~Iwaniec and E.~Kowalski,
\emph{Analytic Number Theory},
American Mathematical Society Colloquium Publications, vol.~53,
American Mathematical Society, Providence, RI, 2004.

\bibitem{littlewood1912}
J.~E.~Littlewood,
\emph{Quelques cons\'equences de l'hypoth\`ese que la fonction
\(\zeta(s)\) de Riemann n'a pas de z\'eros dans le demi-plan
\(\operatorname{Re}s>1/2\)},
C. R. Acad. Sci. Paris \textbf{154} (1912), 263--266.

\bibitem{PascadiPi}
A.~Pascadi,
\emph{Large sieve inequalities for exceptional Maass forms and the
greatest prime factor of \(n^2+1\)},
Forum Math. Pi \textbf{14} (2026), e8.
\href{https://doi.org/10.1017/fmp.2026.10025}{doi:10.1017/fmp.2026.10025}.
Preprint version \href{https://arxiv.org/abs/2404.04239v3}{arXiv:2404.04239v3}.

\bibitem{PascadiFractions}
A.~Pascadi,
\emph{Sums of Kloosterman fractions and applications},
preprint, \href{https://arxiv.org/abs/2611.34917v1}{arXiv:2611.34917v1}, 2026.

\bibitem{Shp}
I.~E.~Shparlinski,
\emph{Bounds of trilinear sums with Kloosterman fractions},
Arch. Math. \textbf{117} (2021), 261--266.
\href{https://doi.org/10.1007/s00013-021-01623-y}{doi:10.1007/s00013-021-01623-y}.

\end{thebibliography}
\end{document}
